The sweep in Table~\ref{tab:sweep} runs at a smaller budget than the headline comparison in
Table~\ref{tab:rl-ce} --- 1{,}500 episodes (150{,}000 transitions) and 3 seeds rather than
4{,}000 episodes (400{,}000 transitions) and 5 seeds --- so its absolute regrets are higher
throughout and are not directly comparable to Table~\ref{tab:rl-ce}. It is used only to
\emph{select} the tuned configuration, which is then re-run at the full budget.

Three findings are worth stating.

\paragraph{Initialisation, not the step size, is the dominant lever.} $Q_0 = 0$ looks neutral but is
strongly \emph{optimistic} in this MDP: every reward is negative, so $V^{*} < 0$ everywhere and a
zero-initialised table over-values every unvisited action. Setting $Q_0 \approx \bar{V}^{*} = -560$
alone moves regret from $12.60\%$ to $5.51\%$ --- more than any other single change, and more than
the step-size exponent contributes.

\paragraph{The step-size schedule matters, in the direction theory predicts.} At a fixed
150,000-transition budget the polynomial schedule $\alpha_n = (1+n)^{-p}$ substantially outperforms
a constant step: constant $\alpha = 0.1$ and $\alpha = 0.5$ give $24.52\%$ and $21.80\%$ regret
against $12.60\%$ for the baseline schedule. Robbins--Monro conditions say a constant $\alpha$ will
not converge asymptotically; that is \emph{not} what this table measures, since every arm here
stops at the same finite budget. What we observe is worse regret at a fixed budget, which is
consistent with non-convergence but does not demonstrate it.
Within the polynomial family, larger $p$ is better here ($p = 1.0$: $8.72\%$; $p = 0.7$: $12.60\%$;
$p = 0.5$: $17.35\%$).

\paragraph{Exploring starts buy coverage, and coverage buys regret.} Removing exploring starts
collapses state--action coverage from $51.6\%$ to $33.7\%$ and degrades regret to $14.86\%$. Since
coverage bounds what any tabular method can know, we report it beside every learning result rather
than only reporting regret.

The $\gamma$ rows are measured against each $\gamma$'s own $V^{*}$ and so are not comparable to the
others; they are included to show that the learning problem gets harder as the effective horizon
lengthens, which is the expected direction.

\begin{table}[h]
\centering
\small
\caption{$Q$-learning hyperparameter sweep. 15 configurations $\times$ 3 seeds at
150,000 transitions each. Baseline is $\alpha_n = (1+n)^{-0.7}$, $\epsilon$ floor $0.05$,
$Q_0 = 0$, exploring starts on.}
\label{tab:sweep}
\begin{tabular}{@{}lrrr@{}}
\toprule
\textbf{Configuration} & \textbf{Mean regret} & \textbf{sd} & \textbf{Coverage} \\
\midrule
baseline                                   & 12.60\% & 1.46 & 51.6\% \\
\midrule
\multicolumn{4}{@{}l}{\itshape Step-size schedule} \\
$\alpha$ power 0.5                         & 17.35\% & 0.93 & 51.6\% \\
$\alpha$ power 1.0                         &  8.72\% & 0.99 & 51.4\% \\
constant $\alpha = 0.1$                    & 24.52\% & 2.72 & 52.2\% \\
constant $\alpha = 0.5$                    & 21.80\% & 1.04 & 52.0\% \\
\midrule
\multicolumn{4}{@{}l}{\itshape Exploration} \\
$\epsilon$ floor 0.0                       & 12.10\% & 0.86 & 51.5\% \\
$\epsilon$ floor 0.2                       & 11.65\% & 0.22 & 51.4\% \\
$\epsilon$ floor 0.5                       & 10.81\% & 1.00 & 50.6\% \\
no exploring starts                        & 14.86\% & 0.37 & \textbf{33.7\%} \\
\midrule
\multicolumn{4}{@{}l}{\itshape Initialisation} \\
neutral $Q_0 = -560 \approx \bar{V}^{*}$   &  \textbf{5.51\%} & 0.36 & 49.8\% \\
pessimistic $Q_0 = -2000$                  &  9.96\% & 0.21 & 45.6\% \\
\midrule
\textbf{tuned: $Q_0 = -560$, $\alpha$ power 1.0} & \textbf{4.78\%} & 0.14 & 49.8\% \\
\midrule
\multicolumn{4}{@{}l}{\itshape Discount (each measured against its own $V^{*}$; not comparable above)} \\
$\gamma = 0.5$                             &  3.41\% & 0.08 & 52.0\% \\
$\gamma = 0.9$                             &  8.80\% & 0.03 & 51.6\% \\
$\gamma = 0.999$                           & 13.92\% & 1.19 & 51.4\% \\
\bottomrule
\end{tabular}
\end{table}
